\chapter{Model and method}
\label{ch:method}

\section{Model}

The angle $\theta$ of the string from the vertical obeys
\begin{equation}
  \frac{\dd^2\theta}{\dd t^2} = -\frac{g}{L} \sin\theta .
  \label{eq:motion}
\end{equation}
For small $\theta$, $\sin\theta \approx \theta$ and \cref{eq:motion} describes simple harmonic motion with period
\begin{equation}
  T = 2\pi \sqrt{L/g} .
  \label{eq:period}
\end{equation}
At an amplitude $\theta_0$ the true period is longer by a factor of about $1 + \theta_0^2/16$~\cite{nelson1986}, which is \num{1.7e-4} at the \qty{0.052}{\radian} used here.

\section{Measurement}

% Apparatus, procedure and how each uncertainty was estimated.
A steel bob hung from a clamped string, released from rest at \qty{3}{\degree}. A light gate timed 20 oscillations at each length, and each period is the mean of five such runs with a standard error of \qty{1}{\milli\second}. Lengths were measured from the clamp to the centre of the bob with a metre rule to \qty{1}{\milli\metre}.
